\section{The corner value: from exact rows to the first correction}
\label{sec:corner}

This section calculates the corner value of the Lebesgue function. Recall that the Lebesgue
function at a point is the sum of the absolute values of all cardinal
polynomials at that point. Its corner value is denoted by
$\Lambda_c(n)$. The calculation has three stages. First we turn the
finite sum into an exact sum of one-dimensional rows. Next we compare
with an equal-mesh problem to identify the limiting constant. Finally
we estimate the exact rows finely enough to find the term of size
$1/n$. Each stage uses a different feature of the same finite sum.

Two elementary ideas organize the calculation. A singular expression
can be easier to sum after its poles, meaning reciprocal singular
terms such as $c/(x-x_0)$, have been written explicitly;
what remains then has bounded derivatives. Also, a displacement of
size $1/n$ in a mesh can produce a constant change in a sum when the
summands are large near a diagonal. We shall see both effects. The
singularities in intermediate trigonometric quotients require care,
but the underlying cardinal polynomials are finite everywhere.

A Cauchy sum here means a sum containing reciprocal distances
$|x-y|^{-1}$. At a corner the absolute cardinal kernel becomes a positive Cauchy sum.
Its two meshes have sizes $n$ and $n+1$. The equal-mesh reference gives
the leading square; the mesh displacement supplies the constant defect
and the first correction. We retain that displacement throughout the proof.
Put $m=n+1$ and, for the rest of this section, write
\[
 \theta_j=j\pi/n,\quad \phi_k=k\pi/m,\qquad
 \varepsilon_j=\begin{cases}1/2,&j=0,n,\\1,&0<j<n,\end{cases}
 \quad \eta_k=\begin{cases}1/2,&k=0,m,\\1,&0<k<m.\end{cases}
\]
Every unspecified constant below is independent of $n\geq2$ and of the
row indices in their stated ranges. Set $A(z)=a(\log z+d)$, with the
constants in \eqref{eq:main-constants}.

Here and below $\log$ is the natural logarithm, $a=2/\pi$ and
$d=\gamma+\log(8/\pi)$. Thus $A(z)^2$ is an explicitly known
quadratic polynomial in $\log z$. The notation $O(E_n)$ means that
the absolute value of the omitted quantity is at most $CE_n$ for
some fixed finite constant $C$. A statement uniform in a row index
means that the same $C$ works for every row in the stated range.
This distinction will matter: an estimate proved with $r$ fixed as
$n$ grows does not automatically control the sum over $1\leq r\leq n$.

The weights $\varepsilon_j$ and $\eta_k$ are endpoint weights.
They implement the trapezoidal convention on the two angular grids:
an interior sample has weight one, and an endpoint sample has weight
one half. Their occurrence is also a consequence of reflection.
An interior angle has two representatives on a full period, whereas
an endpoint is fixed by reflection. The meshes differ by one in their
denominators throughout this section; in particular $\phi_n<\pi$,
while $\phi_m=\pi$.

\subsection{An exact row formula}

The first task is to understand the absolute values in the corner
sum. In its Cauchy denominators the coordinates are
$x_j=\cos\theta_j$ and $y_k=\cos\phi_k$.
These denominators become small when the angles
are close. The useful fact is that their signs, before taking absolute
values, change at a single position in each row. That allows an
absolute row to be reconstructed from its signed sum and its positive
part.

We begin with the exact kernel so that the endpoint contribution is
kept visible. The function $f$ below is used only when its two cosine
arguments are distinct. At an actual collision the original
trigonometric polynomial supplies the value; a quotient formula alone
cannot decide that value.
For $u\ne v$ let $f(u,v)=(1+\cos v)/|\cos u-\cos v|$. We first obtain
\begin{equation}\label{cu:positive}
 \Lambda_c(n)=\frac4{nm}
 \sum_{\substack{0\leq j,k\leq n\\j+k\text{ odd}}}
       \varepsilon_j\eta_k f(\theta_j,\phi_k)+\frac1m.
\end{equation}
To see this directly from \eqref{eq:cardinal-kernel}, define
$D_n(u,v)=\sum_{|p|+|q|\leq n}e^{i(pu+qv)}$. Grouping signs identifies
this with the product Chebyshev sum at a corner.

For clarity, the Chebyshev identity being used is
$T_p(\cos u)=\cos(pu)$. Combining the positive and negative
frequencies gives
\[
 \sum_{p=-P}^{P}e^{ipu}=1+2\sum_{p=1}^{P}\cos(pu).
\]
Applying this grouping in both coordinates shows that the diamond
sum $D_n$ has precisely the multiplicities $\nu_p\nu_q$ in the
kernel $K_n$, with total degree condition $p+q\leq n$.
At $(1,1)$ all the factors $T_p(1)$ equal one. At another corner,
$T_p(-1)=(-1)^p$ changes phases or reflects the node indices.
The separate correction $T_n(z_1)T_n(B_1)$ in the cardinal formula
remains essential after this grouping: it is the term that later
appears as $(-1)^j$. With
$S=(u+v)/2$, $T=(v-u)/2$, finite geometric sums give
\[
 D_n(u,v)=\frac{\sin((n+1)S)\sin((n+1)T)+\sin(nS)\sin(nT)}
                  {\sin S\sin T}.
\]
Indeed $(p,q)\mapsto(p+q,q-p)$ takes the diamond into the square
$[-n,n]^2$ with equal coordinate parity; summing its two parity classes
gives the displayed product.

The change of variables can be checked directly. If
$A=p+q$ and $B=q-p$, then
$p=(A-B)/2$, $q=(A+B)/2$, and
\[
 |p|+|q|=\max(|A|,|B|),\qquad pu+qv=AS+BT.
\]
The inverse indices are integers exactly when $A$ and $B$ have the
same parity. Thus the diamond is the union of an even-even and an
odd-odd rectangular grid in $(A,B)$. For a step-two progression,
\[
 \sum_{h=0}^{M}e^{i(-M+2h)z}
       =\frac{\sin((M+1)z)}{\sin z}.
\]
The two parity classes have progression lengths $n+1$ and $n$.
Multiplying their one-dimensional sums gives the two products of
sines in the formula for $D_n$. This derivation also explains why
the apparent singularities are removable: each sine quotient was
obtained from a finite sum of exponentials. Quotients at zero denominators have their
continuous polynomial values. Reflection of node indices about either
nonnode corner turns the checkerboard odd and leaves its weights unchanged.
The corresponding Lebesgue sum is therefore
$\sum_{j+k\text{ odd}}2\varepsilon_j\eta_k
 |D_n(\theta_j,\phi_k)-(-1)^j|/(nm)$.
Here the mesh identities are
$n\theta_j=j\pi$ and $m\phi_k=k\pi$. For example,
$\sin(m\theta_j)=(-1)^j\sin\theta_j$ and
$\sin(n\phi_k)=(-1)^{k+1}\sin\phi_k$.
Together with $j+k$ odd, these identities reduce the numerator of
$D_n$ to the expression below. The conversion from products of sines
to cosine differences uses
\[
 \cos u-\cos v
       =2\sin\frac{u+v}{2}\sin\frac{v-u}{2}.
\]
Once the correction is subtracted, the numerator simplifies further:
\[
 \frac{2+\cos u+\cos v}{\cos u-\cos v}-1
       =\frac{2(1+\cos v)}{\cos u-\cos v}.
\]
Since $1+\cos v\geq0$ on $[0,\pi]$, taking absolute values now
produces exactly the positive summand $2f(u,v)$.
\[
 D_n(\theta_j,\phi_k)=(-1)^j
       \frac{2+\cos\theta_j+\cos\phi_k}{\cos\theta_j-\cos\phi_k},
 \qquad |D_n(\theta_j,\phi_k)-(-1)^j|=2f(\theta_j,\phi_k).
\]
Coprimality of $n,m$ leaves only the collisions $(0,0)$ and $(n,m)$.
The odd checkerboard excludes the first. At the second,
$D_n(\pi,\pi)=(-1)^n(2n+1)$, so the corrected absolute value $2n$
with weight $1/(2nm)$ contributes $1/m$. The other entries with $k=m$
vanish. This proves \eqref{cu:positive} and the equality of the two
nonnode corner values. The node corners have value one by cardinality.

There are two small points in this reduction worth checking. First,
$\cos u=\cos v$ with $u,v\in[0,\pi]$ implies $u=v$,
because cosine is strictly decreasing there. Hence a collision of
the two meshes requires $jm=kn$. Since consecutive integers are
coprime, $j=0,k=0$ or $j=n,k=m$ are the only possibilities.
Second, the corner $k=m$ is absent from the double sum in
\eqref{cu:positive}. For $j<n$ its numerator is zero and its
contribution vanishes. At $j=n$ the quotient would be $0/0$,
so we evaluate the polynomial instead. Its endpoint weights give
$2\varepsilon_n\eta_m/(nm)=1/(2nm)$; multiplying by the
corrected absolute value $2n$ gives the displayed $1/m$.

A removable singularity is a point at which a formula is undefined
although the function it represents has a finite continuous extension.
The sine-quotient representation of the finite polynomial $D_n$
has this property. The function $f$ is used only at distinct sampled
angles; it has no joint continuous extension at $(\pi,\pi)$.
Neighboring unequal-mesh entries can be large. We keep their Cauchy denominators
because their size creates the logarithmic growth and the limiting
mesh defect.

The sign change within each row makes \eqref{cu:positive} effectively
one dimensional. We call a function analytic at a point when it is
represented there by a convergent power series on a neighborhood.
For an endpoint this neighborhood extends across the endpoint. Introduce
\begin{equation}\label{cu:wB}
 w(x)=\cot(\pi x/2),\qquad
 B(x)=\cot(\pi x)-\frac1{\pi x}+\frac1{\pi(1-x)}.
\end{equation}
Both poles have been removed from $B$; it is analytic across $[0,1]$,
with $B(0)=1/\pi$, $B(1)=-1/\pi$.

The subtraction in $B$ removes a pole at each end. Indeed, as
$x\to0$,
\[
 \cot(\pi x)=\frac1{\pi x}-\frac{\pi x}{3}+O(x^3),
 \qquad \frac1{\pi(1-x)}=\frac1\pi+O(x).
\]
This gives $B(0)=1/\pi$. At the other end put $y=1-x$ and use
$\cot(\pi(1-y))=-\cot(\pi y)$ to obtain $B(1)=-1/\pi$.
The same convergent power series give the asserted analytic
extensions. In particular
$B$ and all derivatives needed below are bounded on the closed
interval. The function $w$ has only the left endpoint pole:
\[
 w(x)=\frac{a}{x}+O(x)\quad(x\downarrow0),\qquad
 w(1-y)=\frac{\pi y}{2}+O(y^3)\quad(y\downarrow0).
\]
Its zero at the right endpoint will compensate for several large
row terms there. For $1\leq r<m$ put
$x=r/m$, $s=m-r$, $\alpha=(1-x)/2$ and
\[
 T_{n,r}=\sum_{\ell=0}^{r-1}\cot\frac{\pi(\ell+\alpha)}n,
 \qquad V_n=\sum_{r=1}^{m-1}w(r/m)
                    \sum_{\ell=0}^{r-1}B((\ell+\alpha)/n).
\]
We use a convergent difference of reciprocal sums to define $\psi$.
For fixed $z>0$, its $k$th summand is
$(z-1)/((k+1)(k+z))=O((k+1)^{-2})$, so the series converges
absolutely. For the moment define $\psi$ by its convergent harmonic series,
\begin{equation}\label{cu:psi}
 \psi(z)=-\gamma+\sum_{k=0}^\infty
          \left(\frac1{k+1}-\frac1{k+z}\right),\qquad z>0,
\end{equation}
and set
\begin{equation}\label{cu:D}
 D_{m,r}=\psi(r+\alpha)+\psi(s-\alpha)-\psi(m-\alpha)-\psi(\alpha),
 \qquad \mathcal H_m=\frac4{\pi m}\sum_{r=1}^{m-1}w(r/m)D_{m,r}.
\end{equation}
We claim the exact identities
\begin{equation}\label{cu:exact}
 \Lambda_c(n)=1-\frac{2n}{m}+\frac4{nm}\sum_{r=1}^{m-1}w(r/m)T_{n,r}
       =-1+\frac2m+\mathcal H_m+\frac{4V_n}{nm}.
\end{equation}
The complete cotangent sum used in this bookkeeping has an elementary
product proof. Up to a nonzero constant independent of $z$,
\[
 \sin(nz)=C_n\prod_{\ell=0}^{n-1}\sin(z+\ell\pi/n).
\]
One way to verify the proportionality is to write each sine in terms
of $e^{iz}$; the product then factors $e^{2inz}-1$ into its $n$
roots. On an interval avoiding zeros, divide the derivative of each
side by the side itself. This operation is called a logarithmic
derivative and gives
\[
 n\cot(nz)=\sum_{\ell=0}^{n-1}\cot(z+\ell\pi/n).
\]
The product identity is global, so this differentiated equality holds
wherever its denominators are nonzero. Differentiation once more yields
\[
 \sum_{\ell=0}^{n-1}\csc^2(z+\ell\pi/n)
                    =n^2\csc^2(nz).
\]
In particular $z=\pi/(2n)$ gives the half-step $n^2$ sum used
for the endpoint row.

For an interior row the sampled indices are $j=1-r+2\ell$.
The two cotangent sums arising from $F_r$ each reduce to
$n\tan(v/2)$. For example, the initial argument of the plus
sum, multiplied by $n$, is
\[
 \frac{nv}{2}+\frac{(1-r)\pi}{2}
      =\frac\pi2-\frac v2,
 \qquad v=\frac{r\pi}{m},
\]
so its cotangent is $\tan(v/2)$. For the minus sum the
corresponding argument differs from $-\pi/2-v/2$ by an
integer multiple of $\pi$, giving the same value. Multiplication
by $w(r/m)/2$ leaves the signed row sum $n$ because
$w(r/m)\tan(v/2)=1$.

Here is the endpoint bookkeeping behind them. Extend a row periodically
and sample its $n$ indices $j\equiv1-r\pmod2$ modulo $2n$. For
$v=r\pi/m$ its signed entries satisfy
\[
 F_r(u):=\frac{1+\cos v}{\cos u-\cos v}
   =\frac{w(r/m)}2\left(\cot\frac{v-u}{2}+\cot\frac{v+u}{2}\right),
 \qquad \sum_jF_r(j\pi/n)=n.
\]
The last equality follows from
$\sum_{\ell=0}^{n-1}\cot(z+\ell\pi/n)=n\cot(nz)$, the logarithmic
derivative of the finite sine product; its right hand cotangent here is
$\tan(v/2)$. The positive samples have representatives
$j=-r+1,-r+3,\ldots,r-1$, since $r-1<nr/m<r$.
Their sum is $P_r=w(r/m)T_{n,r}$, by symmetry and the substitution
$j=r-1-2\ell$. Thus the absolute periodic row is $2P_r-n$.
Folding back to $0\leq j\leq n$ divides by two and produces exactly
$\varepsilon_j$. For $r=0$ the periodic row is $n^2$, since the derivative
of the same cotangent identity gives
$\sum_{\ell=0}^{n-1}\csc^2((\ell+1/2)\pi/n)=n^2$.
Its weight $\eta_0=1/2$, together with the exceptional $1/m$, yields
$\Lambda_c=1/m+n/m+2\sum_{r=1}^{m-1}(2P_r-n)/(nm)$.
This is the first identity in \eqref{cu:exact}. Finally, removing both
poles in each cotangent and using $\psi(z+1)-\psi(z)=1/z$ gives
$T_{n,r}=nD_{m,r}/\pi+\sum_{\ell<r}B((\ell+\alpha)/n)$,
which proves the second identity.

We can spell out both pieces of this row calculation. For a fixed
$r$, the period consists of the $n$ indices of the required parity
modulo $2n$. The evenness and periodicity of cosine pair $j$ with
$2n-j$. An interior index in $\{1,\ldots,n-1\}$ therefore occurs
twice in the periodic sum, while $0$ and $n$, when present in that
parity class, occur once. Consequently
\[
 \frac12\sum_{\substack{j\bmod 2n\\j\equiv1-r\ (2)}}
        |F_r(j\pi/n)|
 =\sum_{\substack{0\leq j\leq n\\j+r\text{ odd}}}
        \varepsilon_j|F_r(\theta_j)|.
\]
This is the precise meaning of folding a periodic row. It covers
both parities of $n$ with the same weights.

The positive indices can be identified directly.
On a symmetric period, $F_r(u)>0$ exactly when $|u|<v$.
For an index this condition is $|j|<nv/\pi=nr/m$.
Since $nr/m=r-r/m$ lies strictly between $r-1$ and $r$,
the parity-compatible indices in that interval are exactly the
listed ones. There are $r$ of them. Substitute $j=r-1-2\ell$:
\[
 \frac{v-j\pi/n}{2}
 =\frac\pi n\left(\ell+\frac{1-r/m}{2}\right).
\]
The other cotangent term in $F_r$ runs through the same $r$
arguments in reverse order. Its prefactor $w/2$ consequently leaves
$wT_{n,r}$ as the positive sum. If $P$ is a positive sum and $Q$
is the magnitude of the negative sum, the signed total is $P-Q=n$,
while the absolute total is $P+Q=2P-n$. This accounts for the
subtraction in \eqref{cu:exact}.

The endpoint $r=0$ can also be evaluated without a limiting row
argument. Here
\[
 F_0(u)=\frac2{\cos u-1}=-\csc^2(u/2).
\]
The parity indices are odd, so its absolute periodic sum is the
$csc^2$ sum stated above. For example, when $n=2$ the two periodic
angles are $\pi/2$ and $3\pi/2$, each giving absolute value $2$.
Their periodic sum is $4=n^2$, and folding divides it by two.
Multiplication by the outer endpoint weight supplies $n/m$ in the
Lebesgue sum. Adding this, the exceptional $1/m$, and the $m-1=n$
interior rows gives
\[
 \frac1m+\frac nm-\frac{2(m-1)}m
       +\frac4{nm}\sum_{r=1}^{m-1}P_r
   =1-\frac{2n}{m}+\frac4{nm}\sum_{r=1}^{m-1}P_r.
\]
Thus even the scalar term in the exact formula records endpoint
information.

Finally, the two-pole decomposition is an exact equality term by
term:
\[
 \cot\frac{\pi(\ell+\alpha)}n
 =\frac n\pi\left(\frac1{\ell+\alpha}
                   -\frac1{n-\ell-\alpha}\right)
      +B((\ell+\alpha)/n).
\]
The recurrence for $\psi$ follows by shifting its convergent series:
$\psi(z+1)-\psi(z)=1/z$. Telescoping gives
\[
 \begin{aligned}
 \sum_{\ell=0}^{r-1}\frac1{\ell+\alpha}
       &=\psi(r+\alpha)-\psi(\alpha),\\
 \sum_{\ell=0}^{r-1}\frac1{n-\ell-\alpha}
       &=\psi(n+1-\alpha)-\psi(n-r+1-\alpha)\\
       &=\psi(m-\alpha)-\psi(s-\alpha).
 \end{aligned}
\]
Their difference is $D_{m,r}$. Thus $\psi$ is simply convenient
notation for the finite reciprocal sums; no special-function theory
is assumed in obtaining the exact identity.

\subsection{The square reference and its mesh defect}

The equal-mesh reference replaces $\phi_k=k\pi/(n+1)$ by
$\theta_k=k\pi/n$. This artificial reference is useful because it
factorizes exactly. The checkerboard condition keeps $j\ne k$,
so there are no diagonal collisions in its Cauchy sum. Its square
structure will determine the coefficient of $\log^2 n$ and the
coefficient of $\log n$ with little calculation. The true problem
retains a finite difference from that square.

The comparison is used here only to identify a limit. It will not
supply an error small enough for the final $1/n$ term. For that term
we shall return to the exact unequal-mesh row formula. Separating
these two jobs makes the later constant identification transparent.
The constant in the answer can be identified before calculating any
first corrections. Define
\[
 L_n=\frac2n\sum_{\substack{1\leq r\leq n\\r\text{ odd}}}
 d_r\csc\frac{r\pi}{2n},\qquad d_r=\begin{cases}1/2,&r=n,\\1,&r<n.\end{cases}
\]
The equal-mesh corner sum is exactly a square:
\begin{equation}\label{cu:square}
 \frac4{n^2}\sum_{\substack{0\leq j,k\leq n\\j+k\text{ odd}}}
 \varepsilon_j\varepsilon_k f(\theta_j,\theta_k)=L_n^2.
\end{equation}
Simultaneous reflection of $j,k$ replaces the numerator by $1-\cos\theta_k$,
so averaging replaces it by one. For $j>k$, put $s=j-k$, $r=j+k$ and
$c_r=\csc(r\pi/(2n))$. Both indices are odd and the denominator is
$2/(c_rc_s)$. The two orders cancel the factor two. For fixed $s<n$,
the values $r=s,2n-s$ have half weight and the intervening values weight
one; their sum after reflection $c_{2n-r}=c_r$ is
$c_s+2\sum_{s<r<n,\ r\text{ odd}}c_r+\boldsymbol1_{n\text{ odd}}c_n$.
The sole pair for $s=n$, when present, has weight $1/4$.
These are precisely the terms in $(\sum_{r\text{ odd}}d_rc_r)^2$,
proving \eqref{cu:square} including its endpoint.

Here is the geometry behind this factorization. After averaging
the simultaneous reflection, the denominator can be written as
\[
 |\cos\theta_j-\cos\theta_k|
 =2\sin\frac{(j+k)\pi}{2n}
       \sin\frac{|j-k|\pi}{2n}.
\]
On the half $j>k$, both sine factors are nonnegative. The variables
$s=j-k$ and $r=j+k$ have odd parity because $j+k$ is odd, and
$j=(r+s)/2$, $k=(r-s)/2$. Thus the new indices run through a
triangular region with $s\leq r\leq2n-s$. The reflection
$r\mapsto2n-r$ reduces this to the triangular part of a square.
The diagonal of that square is represented by the two half-weight
boundary samples, and the off-diagonal entries appear twice.
This is exactly the expansion
\[
 \left(\sum d_rc_r\right)^2
     =\sum d_r^2c_r^2+2\sum_{s<r}d_sd_rc_sc_r.
\]
If $n$ is odd, the final index $r=n$ has $d_n=1/2$, so its
square coefficient is $1/4$. This is why the final endpoint pair
must be retained in the original finite sum.

For a concrete check take $n=2$. The odd reference index is $r=1$,
so $L_2=\csc(\pi/4)=\sqrt2$. The four allowed ordered grid
pairs $(0,1),(1,0),(1,2),(2,1)$ give weighted summands with total
$2$ after the prefactor $4/n^2=1$. Thus the square is $L_2^2=2$.
The example illustrates the endpoint weights in a setting where
every term can be inspected directly.

Pairing the midpoints $y_r=(2r-1)\pi/(4n)$ gives, by
$\cot y+\tan y=2\csc(2y)$,
\[
 \begin{aligned}
 L_n&=\frac1n\sum_{r=1}^n\cot y_r\\
    &=a(2H_{2n}-H_n)+a\log(2/\pi)+O(n^{-2})=A(n)+O(n^{-2}),
 \end{aligned}
\]
Here $H_k=\sum_{j=1}^k j^{-1}$. For the second equality, $\cot y-1/y$
is smooth on $[0,\pi/2]$,
its integral is $\log(2/\pi)$, and the midpoint error is $O(n^{-2})$.
The last equality uses $H_k=\log k+\gamma+1/(2k)+O(k^{-2})$;
the reciprocal terms in $2H_{2n}-H_n$ cancel.

To see how the harmonic sum enters, pair the midpoint angles
$y_r$ and $\pi/2-y_r$. The identity
$\cot y+\tan y=2\csc(2y)$ turns the $n$ cotangents into the
odd cosecants defining $L_n$. When $n$ is odd there is one central
midpoint $y=\pi/4$; the half weight in $L_n$ handles precisely this
unpaired term.

Now separate the pole of cotangent. Since
$y_r=\pi(r-1/2)/(2n)$,
\[
 \frac1n\sum_{r=1}^n\frac1{y_r}
 =a\sum_{r=1}^n\frac1{r-1/2}
 =a(2H_{2n}-H_n).
\]
The remaining function $q(y)=\cot y-1/y$ is smooth at zero.
Its normalized midpoint sum tends to
\[
 \frac2\pi\int_0^{\pi/2}q(y)\,dy
 =a\left[\log\frac{\sin y}{y}\right]_{0}^{\pi/2}
 =a\log(2/\pi).
\]
The midpoint estimate gives an error of order $n^{-2}$ because
$q''$ is bounded on the full closed interval. Substituting the
harmonic expansion gives explicitly
\[
 2H_{2n}-H_n
 =\log n+\gamma+2\log2+O(n^{-2});
\]
the two terms $2/(4n)$ and $-1/(2n)$ cancel. Combining constants,
$2\log2+\log(2/\pi)=\log(8/\pi)$, yields $A(n)$.
In particular
\[
 L_n^2-A(n)^2=(L_n-A(n))(L_n+A(n))
                   =O(\log n/n^2).
\]
The square reference therefore introduces no constant defect and no
term of order $1/n$ on its own.

A small mesh displacement is amplified near $j=k$.
For an offset $|j-k|=\ell$, the angular distance in the reference
is $\ell\pi/n$, so a Cauchy denominator is of size $\ell/n$
away from the angular endpoints. Changing the second angle by
order $1/n$ changes its reciprocal by order $n/\ell^2$.
There are order $n$ pairs of a fixed offset, and the normalization
is order $n^{-2}$. Hence a single fixed offset can leave a nonzero
constant in the limit. The summability of $\ell^{-2}$ suggests
how to control all offsets at once.

This heuristic must also work near $0$ and $\pi$, where cosine
itself has zero derivative. The following exact comparison supplies
that missing uniformity.

We next justify the limiting defect, including the endpoint rows. A
uniform comparison, with $\ell=|j-k|$ odd, is
\begin{equation}\label{cu:deform}
 |f(\theta_j,\phi_k)-f(\theta_j,\theta_k)|
                    \leq\frac{\pi^2n}{2\ell^2}.
\end{equation}
This bound follows from an exact difference. The cases $k=0$ or $j=n$ have zero
difference. Otherwise put $u=\theta_j$, $v_0=\theta_k$, $v_1=\phi_k$,
$\delta=v_0-v_1$, $V=(v_0+v_1)/2$, $S_i=(u+v_i)/2$ and
$Q_i=|u-v_i|/2$. The order of $u,v_i$ is unchanged. Exact subtraction gives
\[
 |f(u,v_1)-f(u,v_0)|=
 \frac{\cos^2(u/2)\sin V\sin(\delta/2)}
      {\sin S_0\sin S_1\sin Q_0\sin Q_1}.
\]
Since $\sin S_i\geq\cos(u/2)\sin(v_i/2)$, the quotient preceding
$\sin(\delta/2)$ and the $Q_i$ factors is at most
$\cot(v_0/2)+\cot(v_1/2)\leq\pi(m-k)/k$.
Here $\sin z\geq2z/\pi$ and $\cos z\leq\pi/2-z$ on $[0,\pi/2]$
suffice. Multiplying by $\sin(\delta/2)\leq\pi k/(2nm)$ gives
$\pi^2(m-k)/(2nm)$. If $k<j$, both $\sin Q_i\geq\ell/n$.
If $k>j$, put $h=n-j$; then
$\sin Q_0\geq\ell/n$, $\sin Q_1\geq[n(\ell-1)+h]/(nm)$ and
$m-k\leq h$. The resulting bound is
$\pi^2nh/[2\ell\{n(\ell-1)+h\}]\leq\pi^2n/(2\ell^2)$,
since $h\leq n$. This proves \eqref{cu:deform}.

The separate case $k>j$ is the place where unequal meshes come
closest to a new collision. Indeed,
\[
 \phi_k-\theta_j
 =\frac\pi{nm}(nk-mj)
 =\frac\pi{nm}\{n(\ell-1)+(n-j)\}.
\]
For $\ell=1$, the braces reduce to $h=n-j$, which may be as
small as one. An estimate replacing this distance by a fixed
multiple of $\ell/n$ would fail. The numerator bound contains
the compensating factor $m-k\leq h$. Keeping that factor until
the division is what gives a uniform result even in the final row.

For completeness, the final algebraic comparison is equivalent to
\[
 h\ell\leq n(\ell-1)+h
       \quad\Longleftrightarrow\quad
       (\ell-1)h\leq(\ell-1)n.
\]
It follows from $h\leq n$, including the equality case $\ell=1$.
When $k<j$, moving $v_0$ toward zero increases its distance from
$u$, so both distance factors already admit the simpler lower bound.
The exact difference and these two cases establish the same
$C n/\ell^2$ estimate everywhere, including the angular endpoints.
This is the estimate needed for exchanging a limit and an infinite
sum of offsets.

Let $\Delta_n$ be $4/(nm)$ times the weighted sum of the differences in
\eqref{cu:deform}. The $k=n$ reference row is zero, so changing its weight
from $\eta_n$ to $\varepsilon_n$ has no effect. Equations
\eqref{cu:positive} and \eqref{cu:square} consequently give
\[
 \Lambda_c(n)-L_n^2=\Delta_n+(1-L_n^2)/m.
\]
The second term tends to zero.

The extra scalar term in this equation comes from normalization,
not from the local mesh displacement. The reference sum is normalized
by $4/n^2$, while the true sum has $4/(nm)$. The latter coefficient
is $n/m$ times the former. Thus the exact relation is
$\Lambda_c=(n/m)L_n^2+\Delta_n+1/m$, which gives the displayed
difference after subtracting $L_n^2$. Since $L_n=O(\log n)$,
$(1-L_n^2)/m=O(\log^2 n/n)$ and hence tends to zero.
For a constant limit this error is harmless. For a first reciprocal
correction it is too large, which is another reason to use the row
formula again later. For fixed odd $\ell$ and $\sigma=\pm1$,
write $k=j+\sigma\ell$, $t=j/n$. On compact subintervals of $(0,1)$,
the cosine difference factors give uniformly
\[
 \frac{f(\theta_j,\phi_k)-f(\theta_j,\theta_k)}n
 \longrightarrow h_{\sigma,\ell}(t)
 =\frac{w(t)}\pi\left(\frac1{|\sigma\ell-t|}-\frac1\ell\right),
\]
because $\phi_k-\theta_j=\pi(\sigma\ell-t)/m$.

For the local limit, keep $\ell$ fixed and let $j/n=t$ stay in
a compact subinterval $[\delta,1-\delta]$. On this interval
$\sin(\pi t)$ is bounded away from zero. Taylor's formula in the
second angle then gives
\[
 |\cos\theta_j-\cos\phi_k|
  =\sin(\pi t)\,\frac{\pi|\sigma\ell-t|}{m}
       +O(m^{-2}),\qquad
 1+\cos\phi_k=1+\cos(\pi t)+O(m^{-1}).
\]
The ratio $(1+\cos(\pi t))/\sin(\pi t)$ is $w(t)$.
For the reference angle, the same calculation replaces
$|\sigma\ell-t|$ by $\ell$ and $m$ by $n$.
After division by $n$, the difference is the stated function
$h_{\sigma,\ell}$. This argument is uniform on the compact
interval, and the limiting factors $|\sigma\ell-t|$ remain bounded away
from zero there for each fixed odd $\ell$.

The normalization explains the factor in the final integral. A fixed
orientation contributes
\[
 \frac4{nm}\sum_j
   \{f(\theta_j,\phi_{j+\sigma\ell})
             -f(\theta_j,\theta_{j+\sigma\ell})\}
 \longrightarrow4\int_0^1h_{\sigma,\ell}(t)\,dt.
\]
The endpoint weights alter only finitely many entries for this fixed
offset, and their normalized contributions vanish. The two
orientations are both present for every positive odd offset. Bound
\eqref{cu:deform} controls each normalized summand by $C/\ell^2$.
For fixed $\ell$, the rows with $t$ outside $[\delta,1-\delta]$ therefore
contribute at most $C(\delta+1/n)/\ell^2$; the same bound controls the
omitted limiting integrals. Interior Riemann sums, followed by
$\delta\downarrow0$, give the limit for that offset. The tails over
$\ell>M$ are uniformly bounded by $C\sum_{\ell>M}\ell^{-2}$, so they
may then be removed. Pairing orientations leaves
$1/(\ell-t)+1/(\ell+t)-2/\ell=2t^2/[\ell(\ell^2-t^2)]$.
Thus, on the full integer sequence,
\begin{equation}\label{cu:kappa}
 \Lambda_c(n)-A(n)^2\longrightarrow\kappa
\end{equation}
with exactly the integral \eqref{eq:main-kappa}. Its terms with $\ell\geq3$
converge uniformly up to both endpoints; for $\ell=1$ the zero of $w$ at
one cancels the pole. At zero $t^2w(t)\to0$. This also proves the claimed
continuous extension of the integrand.

The passage from local limits to the complete sum has a definite
order. First fix an offset cutoff $M$ and an endpoint width
$\delta>0$. On $[\delta,1-\delta]$ ordinary Riemann sums converge
for the finitely many offsets $\ell\leq M$. The uniform estimate
\eqref{cu:deform} bounds the omitted endpoint strips by a constant
multiple of $\delta+1/n$ for each such offset. We may therefore
let $n$ tend to infinity and then let $\delta$ tend to zero.
Finally the same estimate bounds the offset tail, before and after
the limit, by $C\sum_{\ell>M}\ell^{-2}$, which tends to zero.
This three-step procedure proves convergence of the entire sequence;
it does not rely on separate even and odd subsequences.

Pairing orientations is an algebraic cancellation with a useful
interpretation. Each orientation has a displacement of opposite sign,
so their linear effects cancel. Their combined denominator satisfies
\[
 \frac1{\ell-t}+\frac1{\ell+t}-\frac2\ell
 =\frac{2t^2}{\ell(\ell^2-t^2)}.
\]
Multiplying by $4w(t)/\pi$ gives precisely the integrand and
coefficient in $\kappa$. Thus $\kappa$ measures the surviving
coupled effect of the two meshes. It is not an arbitrary additive
constant fitted after the leading square is found.

The endpoint cancellation can be checked numerically in scale,
without evaluating the integral. For $\ell=1$, as $t\uparrow1$,
\[
 \frac{t^2w(t)}{1-t^2}
   =\frac{t^2\tan(\pi(1-t)/2)}{(1-t)(1+t)}
       \longrightarrow\frac\pi4.
\]
All terms with $\ell\geq3$ tend to zero there after multiplication
by $w(t)$. Hence the full integrand, including the prefactor
$8/\pi$, tends to $2$ at the right endpoint. At zero,
$t^2w(t)=at+O(t^3)$ and the offset series is bounded, so the
integrand tends to zero. The terms are nonnegative in the open
interval; in particular the mesh defect is a positive constant.
These facts establish finiteness and continuity without requiring
a numerical value of $\kappa$.

\subsection{Uniform summation and the first correction}\label{sec:asym}
To obtain the first correction, we return to \eqref{cu:exact}, keeping
the harmonic and smooth parts together.
One elementary quadrature rule will control all remainders.

Quadrature means approximating an integral by a weighted finite sum
of function values. The familiar trapezoidal rule has error
$O(m^{-2})$ when a function has a bounded second derivative.
Some functions below have the milder endpoint behavior $x\log x$
or $(1-x)\log(1-x)$. They are continuous, but their second
derivatives behave like $1/x$ or $1/(1-x)$. Summing the local
trapezoidal errors then produces one logarithm, giving
$O(\log m/m^2)$. The next lemma states exactly this version.

The endpoint term in the lemma also fixes a convention. The sum is
over interior points only. The full trapezoidal approximation would
add $(g(0)+g(1))/(2m)$ to that sum. Solving for the interior sum
therefore gives a minus sign in front of the endpoint correction.
That sign will enter the coefficient of $1/n$.

\begin{lemma}[Quadrature with endpoint logarithms]\label{cu:quadrature}
Suppose $g$ is continuous on $[0,1]$, twice continuously differentiable
inside, with $|g''(x)|\leq C(x^{-1}+(1-x)^{-1})$ there, and, for
$0<x\leq1/2$,
\[
 |g(x)-g(0)|+|g(1-x)-g(1)|\leq Cx(1+|\log x|).
\]
Then
\[
 \frac1m\sum_{r=1}^{m-1}g(r/m)=\int_0^1g(x)\,dx
       -\frac{g(0)+g(1)}{2m}+O(\log m/m^2).
\]
For $g\in C^2[0,1]$ the error is $O(m^{-2})$; its composite midpoint
rule has the same order.
\end{lemma}
\begin{proof}
To track the sign and scale of the cell error, set
$K(t)=(t-u)(u+h-t)/2$. This function vanishes at the two ends of
the cell, has $K'(u)=h/2$, $K'(u+h)=-h/2$, and $K''=-1$.
Integrating $Kg''$ by parts twice gives the formula below.
The factor $K$ has size at most $h^2/8$, which makes it especially
well suited to integrable bounds for $g''$. On a cell $[u,u+h]$
we therefore have
\[
 \frac h2\{g(u)+g(u+h)\}-\int_u^{u+h}g
 =\frac12\int_u^{u+h}(t-u)(u+h-t)g''(t)\,dt.
\]
For $h=1/m$ the two endpoint cells cost $O(h^2(1+|\log h|))$ by
the assumed moduli. On the others the kernel is at most $h^2/8$,
and $\int_h^{1-h}|g''|=O(\log m)$. For bounded second derivative,
the local error is at most $h^3\|g''\|_\infty/12$. Taylor's formula
at the midpoint gives the corresponding bound $h^3\|g''\|_\infty/24$.
\end{proof}

The endpoint-cell estimate does not require evaluating $g'$ or
$g''$ at an endpoint. On the first cell, subtract the constant
$g(0)$. Both its integral and its trapezoidal sum cancel exactly.
The remaining absolute values are bounded by
\[
 C\left(h\,h(1+|\log h|)
       +\int_0^h x(1+|\log x|)\,dx\right)
       =O(h^2(1+|\log h|)).
\]
The last cell is identical after replacing $x$ by $1-x$.
On the interior cells,
\[
 \int_h^{1-h}\left(\frac1x+\frac1{1-x}\right)dx
                 =O(\log m).
\]
Multiplying this by the kernel bound proves the asserted total error.
The model function $g(x)=x\log x$, extended by $g(0)=0$, illustrates
the assumptions: $g''(x)=1/x$ and the endpoint variation has exactly
the stated logarithmic size. Thus the lemma allows the actual
functions needed below, without assuming bounded endpoint derivatives.

We also record the elementary estimates used below.

The function $\psi$ enters the proof through finite reciprocal
sums, but its derivatives give an efficient way to control the
changes in their arguments. The defining series is differentiable
on every compact interval $[c,C]\subset(0,\infty)$ because the
series of derivatives is dominated by a convergent reciprocal-square
series. Differentiating again is justified by the same argument
with reciprocal cubes. Such convergence is called local uniform
convergence: it is uniform whenever the variable stays a positive
distance from the pole at zero.

For example, integral comparison gives
\[
 \sum_{k=0}^\infty(z+k)^{-3}
 \leq z^{-3}+\int_0^\infty(z+t)^{-3}\,dt
 =z^{-3}+\frac1{2z^2},
\]
which proves the stated bound on $\psi''$. This explicit bound
will turn Taylor's formula into an estimate valid even in the first
and last rows. Differentiating
\eqref{cu:psi} on compact subsets of $z>0$ gives
\[
 \psi'(z)=\sum_{k\geq0}(z+k)^{-2},\quad
 \psi''(z)=-2\sum_{k\geq0}(z+k)^{-3},\quad
 |\psi''(z)|\leq z^{-2}+2z^{-3}.
\]
The unit-cell trapezoid identity applied to $(z+t)^{-1}$ and integral
comparison for $(z+t)^{-2}$ imply, uniformly for $z\geq1$,
\begin{equation}\label{cu:psi-est}
 \psi(z)=\log z-\frac1{2z}+O(z^{-2}),\qquad
 \psi'(z)=z^{-1}+O(z^{-2}).
\end{equation}
Indeed $\psi(z)=\lim_{K\to\infty}(\log K-\sum_{k<K}(z+k)^{-1})$,
and the trapezoid error is bounded by $\tfrac18\int_0^\infty
2(z+t)^{-3}dt$. This also proves the harmonic estimate used above.
The same series gives
$\psi(1/2)=-\gamma-2\log2$, $\psi'(1/2)=\pi^2/2$.
For the latter value, the classical paired partial fraction
$\pi\cot(\pi z)=z^{-1}+\sum_{k\geq1}((z-k)^{-1}+(z+k)^{-1})$
gives $\sum k^{-2}=\pi^2/6$ by its linear term at zero; removing even
indices gives $4\sum_{k\text{ odd}}k^{-2}=\pi^2/2$.
One short derivation of the partial fraction is the cosine series of
$\cos(zs)$ on $[0,\pi]$: its constant coefficient is
$2\sin(\pi z)/(\pi z)$ and its $k$th coefficient is
$2(-1)^kz\sin(\pi z)/(\pi(z^2-k^2))$. The coefficients are absolutely
summable; evaluating at $s=\pi$ gives the identity. All differentiated
series used here converge locally uniformly away from their poles.

For completeness, a cosine-series coefficient of a function $h$
on $[0,\pi]$ means
$a_k=(2/\pi)\int_0^\pi h(s)\cos(ks)\,ds$, with constant
term $a_0/2$. For $h(s)=\cos(zs)$, product to sum and integration
give precisely the coefficients above. For noninteger $z$ their
absolute values are $O(k^{-2})$. The series therefore converges
absolutely and uniformly; the Fourier convergence theorem for a
continuous piecewise smooth even periodic function identifies it
with $h$, also at the endpoints. Evaluating at $s=\pi$ gives
\[
 \pi\cot(\pi z)=\frac1z+
             2z\sum_{k=1}^{\infty}\frac1{z^2-k^2}.
\]
Near zero the right side is
$1/z-2z\sum_{k\geq1}k^{-2}+O(z^3)$; the uniform bound for
$|z|\leq1/2$ justifies taking this coefficient term by term.
The elementary sine and cosine expansions give
$\pi\cot(\pi z)=1/z-\pi^2z/3+O(z^3)$ on the left.
Equating coefficients yields $\sum k^{-2}=\pi^2/6$.
This supplies the one reciprocal-square value used in the proof.

The large-$z$ expansion can also be viewed directly through
quadrature. The series defining $\psi$ is equivalent to
\[
 \psi(z)=\lim_{K\to\infty}
       \left(\log K-\sum_{k=0}^{K-1}\frac1{z+k}\right),
\]
because $H_K-\log K\to\gamma$. Applying trapezoids to
$f(t)=1/(z+t)$ gives
\[
 \sum_{k=0}^{K-1}f(k)
 =\int_0^K f(t)\,dt+\frac{f(0)-f(K)}2+O(z^{-2}),
\]
where the bound is independent of $K$. Since
$\int_0^K f=\log(z+K)-\log z$, passage to the limit yields
$\log z-1/(2z)+O(z^{-2})$ for $\psi(z)$.
For $\psi'$, the integral of $(z+t)^{-2}$ is $1/z$, and the
difference between the series and the integral is $O(z^{-2})$.
These are the two estimates in \eqref{cu:psi-est}.

At half integers, the same formulas have simple exact constants.
The recurrence gives
\[
 \psi(R+1/2)=\psi(1/2)+2H_{2R}-H_R.
\]
Comparing its large-$R$ limit with $\log R$ shows
$\psi(1/2)=-\gamma-2\log2$. Also
\[
 \psi'(1/2)=4\sum_{k=0}^\infty\frac1{(2k+1)^2}
 =4\left(1-\frac14\right)\sum_{k=1}^\infty\frac1{k^2}
 =\frac{\pi^2}{2}.
\]
The partial-fraction argument above supplies the reciprocal-square
sum, so both half-integer values follow from the series and elementary
Fourier coefficients. No table of special-function values is needed.

Define the finite, endpoint-regular integrals
\[
 I_1=\int_0^1\left(\frac{w(x)}{1-x}-\frac ax\right)dx,
 \qquad I_2=\int_0^1xw(x)\,dx,\qquad I_B=\int_0^1xw(x)B(x)\,dx.
\]
The next proposition contains the only aggregate calculation needed.
\begin{proposition}\label{cu:aggregate}
There are constants $C_H,C_B$, independent of the degree, such that
\begin{align}
 \mathcal H_m&=A(m)^2+C_H+
       \frac{-a^2\log m+\beta_H}{m}+O(\log m/m^2),\label{cu:H}\\
 \frac{4V_n}{nm}&=C_B+\frac{\beta_B}{n}+O(n^{-2}),\label{cu:V}\\
 \beta_H&=-2-a^2(\gamma+2\log2)-a(I_1+I_2),\notag\\
 C_B&=4\int_0^1w(t)\int_0^tB(x)\,dx\,dt,
 \qquad \beta_B=-a^2+2I_B+aI_2.\notag
\end{align}
\end{proposition}

The two constants $C_H$ and $C_B$ collect terms independent of the
degree. We will identify their required combination by the already
proved limit $\kappa$, so it is unnecessary to evaluate them
separately. Their reciprocal corrections are different: the harmonic
part has $(-a^2\log m+\beta_H)/m$, while the smooth part has
$\beta_B/n$. The smooth part therefore contributes to the first
correction even though its summands have no poles. The proof keeps
that term and the endpoint sequence from the harmonic part until
they can be combined exactly.
\begin{proof}
We now derive the harmonic contribution in several steps. The
important point is to choose expansion centers that remain useful
when $r$ or $s=m-r$ is small. Replacing every occurrence of $\psi$
by a large-argument formula would be invalid at those ends. Instead,
we retain $\psi(\alpha)$ exactly, center the other arguments at
$r+1/2$, $s$, and $m$, and isolate two convergent sequences from
the small-$r$ end.

First center the harmonic row at its exact integer and half integer
arguments, before expanding. Taylor's formula gives, uniformly for
$r+s=m$, $x=r/m$, $\alpha=(1-x)/2$,
\begin{equation}\label{cu:center}
 \begin{split}
 D_{m,r}={}&\psi(r+1/2)+\psi(s)-\psi(\alpha)-\psi(m)\\
 &-\frac{r\psi'(r+1/2)+s\psi'(s)-s\psi'(m)}{2m}+O(m^{-2}).
 \end{split}
\end{equation}
In fact the three shifts are respectively $r/(2m),s/(2m),s/(2m)$.
For a shift $h$ at $z$ the error is at most
$h^2\{(z-h)^{-2}+2(z-h)^{-3}\}/2$.
The lower arguments are at least $r$, $3s/4$, $3m/4$, respectively;
hence each error is $O(m^{-2})$, even for $r=1$ or $s=1$.

More explicitly, $\alpha=s/(2m)$, so the exact arguments are
\[
 r+\alpha=r+\frac12-\frac{r}{2m},\qquad
 s-\alpha=s-\frac{s}{2m},\qquad
 m-\alpha=m-\frac{s}{2m}.
\]
For the first argument the shift is proportional to $r$, for the
second to $s$. Taylor's second-order remainder is therefore small
uniformly: the square of the shift contains the same square of
the row index that occurs in the denominator of $\psi''$.
For example,
\[
 \left(\frac{r}{2m}\right)^2
       \left(\frac1{r^2}+\frac2{r^3}\right)
 \leq\frac{3}{4m^2}\qquad(r\geq1).
\]
The corresponding estimates at $s$ and $m$ differ only by fixed
factors from the lower-argument bounds in the text.

There is a sign to keep in the third Taylor expansion.
The original term is $-\psi(m-s/(2m))$, so its first derivative
contribution is positive. This explains the subtraction of
$s\psi'(m)$ inside the bracket in \eqref{cu:center}. Keeping
this sign is necessary for the later constant cancellation.
Since $w(r/m)\leq am/r$, its weighted normalized error is
$O(\log m/m^2)$.

Two summable endpoint sequences must still be retained. Put
$e_r=\psi(r+1/2)-\log r$, $u_r=r\psi'(r+1/2)-1$.
The midpoint rule for $t^{-2}$ on $[r,\infty)$ gives
$\psi'(r+1/2)=1/r+O(r^{-3})$: the local errors sum as
$O(\sum_{k\geq0}(r+k)^{-4})$. Integration from infinity, using
\eqref{cu:psi-est}, then gives $e_r=O(r^{-2})$ as well as
$u_r=O(r^{-2})$. Thus $S_e=\sum e_r/r$ and $S_u=\sum u_r/r$ converge.

Here the improved $O(r^{-2})$ errors are stronger than the generic
large-argument estimates in \eqref{cu:psi-est}. They come from
centering at a half integer. The values
$\psi'(r+1/2)=\sum_{k\geq0}(r+k+1/2)^{-2}$ are midpoint
samples of $t^{-2}$ on $[r,\infty)$. The integral is $1/r$,
and the sum of the midpoint errors is $O(r^{-3})$.
Consequently
\[
 \frac{d}{dr}\{\psi(r+1/2)-\log r\}=O(r^{-3})
\]
when $r$ is regarded temporarily as a real variable. The expression
in braces tends to zero by \eqref{cu:psi-est}, so integration from
$r$ to infinity gives $e_r=O(r^{-2})$. Multiplying the derivative
estimate by $r$ gives $u_r=O(r^{-2})$.

These two sequences encode finite endpoint effects. Although each
term decays quickly, the factor $w(r/m)\sim am/r$ makes their
weighted sums tend to nonzero constants. The sequence $e_r$
contributes to the constant term, and $u_r$, which is already
multiplied by $1/m$, contributes to the first correction. Discarding
all fixed-index terms as negligible would lose these contributions.
Finite recurrence summation gives
\[
 \sum_{r=1}^R\psi'(r+1/2)
 =(R+1/2)\psi'(R+1/2)+\psi(R+1/2)
               -\psi(1/2)-\tfrac12\psi'(1/2).
\]
This follows by summing
$\psi'(r+1/2)=\psi'(1/2)-\sum_{k<r}(k+1/2)^{-2}$ and writing
$k=(k+1/2)-1/2$ in the finite triangle. Subtracting $H_R$ and passing
to the limit gives $S_u=1+2\log2-\pi^2/4$.

The finite recurrence identity can be verified without rearranging
an infinite double sum. Write $p_k=\psi'(k+1/2)$. Differentiating
the recurrence for $\psi$ gives
$p_k-p_{k+1}=(k+1/2)^{-2}$. An immediate finite telescoping identity is
\[
 \sum_{r=1}^R p_r
 =(R+1/2)p_R+\psi(R+1/2)
       -\psi(1/2)-\tfrac12p_0.
\]
For example, its difference between $R$ and $R-1$ equals $p_R$,
and at $R=1$ it follows directly from the recurrence.
Since $u_r/r=p_r-1/r$, subtracting $H_R$ from the left gives
the partial sum of $S_u$. On the right,
$(R+1/2)p_R\to1$ and
$\psi(R+1/2)-H_R\to-\gamma$. The half-integer values then give
\[
 S_u=1-\gamma-(-\gamma-2\log2)-\frac{\pi^2}{4}
       =1+2\log2-\frac{\pi^2}{4}.
\]
This calculation retains the finite derivative contribution rather
than hiding it in an unspecified error constant.

Let $d_0=\gamma+2\log2$, $q_0=\pi^2/4-1$ and
$F(x)=\log x+\log(1-x)-\psi((1-x)/2)$. Substitution of
\eqref{cu:psi-est} into \eqref{cu:center}, keeping $e_r,u_r$ intact, yields
\begin{equation}\label{cu:row-expand}
 D_{m,r}=\log m+F(x)+e_r-\frac1{2m(1-x)}-\frac{x+u_r}{2m}+Q_{m,r},
 \quad |Q_{m,r}|\leq C\{s^{-2}+(ms)^{-1}+m^{-2}\}.
\end{equation}
Indeed $\psi(s)-\psi(m)$ contributes
$\log(s/m)-1/(2s)+1/(2m)$, whereas the derivative bracket is
$1+x+u_r+O(1/s+1/m)$; its constant cancels $1/(2m)$.

To see each term of the row expansion, start with
$\psi(r+1/2)=\log r+e_r$ and
$r\psi'(r+1/2)=1+u_r$. The two large integer arguments give
\[
 \begin{aligned}
 \psi(s)-\psi(m)
   &=\log(1-x)-\frac1{2s}+\frac1{2m}
                         +O(s^{-2}+m^{-2}),\\
 s\psi'(s)&=1+O(s^{-1}),\\
 s\psi'(m)&=1-x+O(m^{-1}).
 \end{aligned}
\]
Hence the derivative bracket in \eqref{cu:center} is
$1+x+u_r+O(s^{-1}+m^{-1})$. Its $-1/(2m)$ cancels the
$+1/(2m)$ from the integer-argument difference. Also
$\log r=\log m+\log x$. These substitutions leave precisely
$F(x)$ and the two displayed reciprocal corrections in
\eqref{cu:row-expand}.

The error bound is deliberately allowed to be large when $s$ is
small. It includes $s^{-2}$, which is order one in the last row.
This is acceptable because $w(r/m)$ vanishes at that end.
A uniform aggregate estimate uses the product of the row error and
its actual weight, rather than demanding that the row error alone
be small everywhere.
To sum the remainder, use $w(r/m)\leq am/r$ on the first half and
$w(r/m)\leq Cs/m$ on $r\geq m/2$. They give respectively bounds
$C/(m^2r)$ and $C/(m^2s)$ after normalization by $1/m$.
Consequently $m^{-1}\sum w|Q|=O(\log m/m^2)$.

Here is the split in detail. If $r\leq m/2$, then $s\geq m/2$,
so $|Q_{m,r}|\leq C/m^2$ and
\[
 \frac1m w(r/m)|Q_{m,r}|\leq\frac{C}{m^2r}.
\]
If $r\geq m/2$, then $w(r/m)\leq Cs/m$ and
\[
 \frac1m w(r/m)|Q_{m,r}|
 \leq \frac{Cs}{m^2}
       \left(\frac1{s^2}+\frac1{ms}+\frac1{m^2}\right)
 \leq\frac{C}{m^2s}.
\]
Summing either bound produces a harmonic sum, at most a constant
multiple of $1+\log m$. Thus the endpoint rows are included in
the error estimate. This is a concrete example of uniform
summation: a row estimate is chosen together with the weights that
will actually be used when summing it.

Write $v(x)=w(x)-a/x$, smoothly extended by $v(0)=0$; $|v(x)|\leq Cx$.
The sequence bounds and their $O(m^{-2})$ tails show
$m^{-1}\sum_{r<m}w(r/m)e_r=aS_e+O(\log m/m^2)$,
and the same assertion holds with $u_r,S_u$. Each $v$ contribution is
bounded by $Cm^{-2}\sum_{r<m}r^{-1}$.

For instance, decompose $w(r/m)=am/r+v(r/m)$. The main part gives
\[
 \frac1m\sum_{r<m}\frac{am}{r}e_r
   =aS_e-a\sum_{r\geq m}\frac{e_r}{r}
   =aS_e+O(m^{-2}),
\]
because $e_r/r=O(r^{-3})$. The smooth remainder obeys
\[
 \frac1m\sum_{r<m}|v(r/m)e_r|
 \leq\frac{C}{m^2}\sum_{r<m}\frac1r
       =O(\log m/m^2).
\]
The identical argument works for $u_r$. This explains why the
infinite constants $S_e$ and $S_u$ may replace their finite sums
at exactly the required precision.

The two ordinary quadrature functions are regular at both endpoints.
Near zero,
\[
 \frac{w(x)}{1-x}-\frac ax
   =a+O(x),\qquad xw(x)=a+O(x^2).
\]
Near one, the zero of $w$ cancels the denominator $1-x$.
Thus their trapezoidal errors have no endpoint logarithm.
The only harmonic singularity in the first sum below is the
explicit $a/x$ part, whose grid sum is $aH_{m-1}$. Ordinary trapezoid estimates for
the smooth functions $w/(1-x)-a/x$ and $xw$ also give
\[
 \frac1m\sum_{r<m}\frac{w(r/m)}{1-r/m}
     =a\log m+a\gamma+I_1+O(m^{-1}),\qquad
 \frac1m\sum_{r<m}(r/m)w(r/m)=I_2+O(m^{-1}).
\]
Thus, for $M_m=m^{-1}\sum_{r<m}w(r/m)\{\log m+F(r/m)\}$,
\begin{equation}\label{cu:H-to-M}
 \mathcal H_m=2a\left\{M_m+aS_e-
    \frac{a\log m+a\gamma+I_1+I_2+aS_u}{2m}\right\}
                      +O(\log m/m^2).
\end{equation}

The quantity $M_m$ contains the logarithmic square. Its singular
part is explicit enough to sum with harmonic sums; its remainder
is continuous but has logarithmic derivatives at the endpoints.
The aim is to isolate these two pieces before applying quadrature.
This avoids using an integral approximation directly on
$\log x/x$, whose integral diverges at zero.

It remains to compute the endpoint correction of $M_m$. Subtract its
explicit singularity:
$R(x)=w(x)F(x)-a(\log x+d_0)/x$.
Near zero $J(x)=\log(1-x)-\psi((1-x)/2)$ is analytic, with
$J(0)=d_0$, $J'(0)=q_0$, and
\[
 R(x)=a\frac{J(x)-d_0}{x}+v(x)\{\log x+J(x)\}.
\]
Near one, write $y=1-x$ and use
$\psi(y/2)=-2/y+\psi(1+y/2)$, $w(1-y)=\tan(\pi y/2)$.
These analytic factors prove
$R(0)=aq_0$, $R(1)=\pi-ad_0$ and exactly the derivative and endpoint
modulus bounds of Lemma~\ref{cu:quadrature}. Its conclusion applies.

We can verify the endpoint hypotheses explicitly. At zero,
write $J(x)=d_0+q_0x+O(x^2)$ and $v(x)=O(x)$.
The first term of $R$ is analytic, while the second contains
only a possible $x\log x$ term. Hence
\[
 R(x)=aq_0+O(x(1+|\log x|)),\qquad
 |R''(x)|\leq C/x
\]
near zero. To compute $q_0$, differentiate $J$:
\[
 J'(x)=-\frac1{1-x}+\frac12\psi'((1-x)/2),
 \qquad J'(0)=-1+\frac{\pi^2}{4}.
\]
At the other endpoint put $y=1-x$. The recurrence gives
\[
 F(1-y)=\log(1-y)+\log y+\frac2y-\psi(1+y/2).
\]
Multiplication by $w(1-y)=\tan(\pi y/2)$ makes the pole
$2/y$ tend to $\pi$. The only nonsmooth surviving factor is
of size $y\log y$. The explicitly subtracted term tends to
$ad_0$, so
\[
 R(1-y)=\pi-ad_0+O(y(1+|\log y|)),\qquad
 |R''(1-y)|\leq C/y.
\]
On any closed interior interval the derivatives are bounded.
Combining the three regions gives precisely the hypotheses of
Lemma~\ref{cu:quadrature} on all of $(0,1)$. Its correction is
therefore $-(aq_0+\pi-ad_0)/(2m)$.

The integral $\int v$ is also elementary. With a lower cutoff
$\epsilon>0$,
\[
 \int_\epsilon^1\left(w(x)-\frac ax\right)dx
 =\left[a\log\sin(\pi x/2)-a\log x\right]_\epsilon^1.
\]
As $\epsilon\downarrow0$, the ratio
$\sin(\pi\epsilon/2)/\epsilon$ tends to $\pi/2$,
so the integral is $a\log(2/\pi)$. This value and
$v(0)=0$, $v(1)=-a$ determine the smooth weight sum.
Also $v(1)=-a$, $\int_0^1v=a\log(2/\pi)$, so smooth quadrature gives
\begin{equation}\label{cu:w-sum}
 \frac1m\sum_{r<m}w(r/m)=a\{\log m+\gamma+\log(2/\pi)\}+O(m^{-2}).
\end{equation}
The $-a/(2m)$ from $aH_{m-1}$ cancels the endpoint correction for $v$.

Indeed the decomposition gives
\[
 \frac1m\sum_{r<m}w(r/m)
 =aH_{m-1}+\frac1m\sum_{r<m}v(r/m).
\]
The first term has reciprocal coefficient $-a/(2m)$.
For the second, trapezoids give
$\int v-(v(0)+v(1))/(2m)+O(m^{-2})
=a\log(2/\pi)+a/(2m)+O(m^{-2})$.
The sum has no reciprocal term. Consequently multiplying this
weight sum by $\log m$ introduces an error of only
$O(\log m/m^2)$, which is the accuracy needed for $M_m$.
Finally, unit-cell trapezoids for $\log t/t$ give a real constant $C_K$ with
\[
 H_{m-1}=\log m+\gamma-\frac1{2m}+O(m^{-2}),\qquad
 \sum_{r<m}\frac{\log r}{r}=\tfrac12\log^2m+C_K
                   -\frac{\log m}{2m}+O(\log m/m^2).
\]
For the second formula the local errors have a convergent sum because
$(\log t/t)''=(2\log t-3)/t^3$; the tail is $O(\log m/m^2)$.
Express $M_m$ as $\log m$ times \eqref{cu:w-sum}, plus
$a\sum_{r<m}(\log r)/r+a(d_0-\log m)H_{m-1}$, plus $m^{-1}\sum R$.
For an explicit check of the reciprocal-term cancellation, abbreviate
$L=\log m$ and $c_v=\log(2/\pi)$. The three singular pieces
of $M_m$ expand as
\[
 \begin{aligned}
 L\frac1m\sum_{r<m}w(r/m)
    &=aL^2+a(\gamma+c_v)L+O(L/m^2),\\
 a\sum_{r<m}\frac{\log r}{r}
    &=\frac a2L^2+aC_K-\frac{aL}{2m}+O(L/m^2),\\
 a(d_0-L)H_{m-1}
    &=-aL^2+a(d_0-\gamma)L+ad_0\gamma
         -\frac{ad_0}{2m}+\frac{aL}{2m}+O(L/m^2).
 \end{aligned}
\]
The $L/m$ terms cancel. The $L$ coefficient is
$a(c_v+d_0)=ad$. The remaining quadrature piece has reciprocal
coefficient $-(aq_0+\pi-ad_0)/(2m)$, which cancels the
$-ad_0/(2m)$ above except for the indicated $\pi+aq_0$.
All constant terms, including $\int R$, may now be collected in
one degree-independent number $C_M$. Thus the resulting formula is
\[
 M_m=\tfrac a2\log^2m+ad\log m+C_M
          -\frac{\pi+aq_0}{2m}+O(\log m/m^2)
\]
for a constant $C_M$. Inserting this into \eqref{cu:H-to-M} and using
$q_0+S_u=2\log2$ proves \eqref{cu:H} and its stated $\beta_H$;
all degree-independent terms are absorbed into $C_H$.

The last simplification can be checked independently of those
constant terms. Multiplying the reciprocal coefficient of $M_m$
by $2a$ gives $-a\pi-a^2q_0$. The bracket in
\eqref{cu:H-to-M} contributes
$-a^2\log m-a^2\gamma-a(I_1+I_2)-a^2S_u$.
Since $a\pi=2$ and
\[
 q_0+S_u
   =\left(\frac{\pi^2}{4}-1\right)
      +\left(1+2\log2-\frac{\pi^2}{4}\right)
   =2\log2,
\]
the constant reciprocal coefficient is exactly
$-2-a^2(\gamma+2\log2)-a(I_1+I_2)=\beta_H$.
The leading terms are $a^2\log^2m+2a^2d\log m$;
adding or subtracting the fixed quantity $a^2d^2$ only changes
$C_H$, so they may be written as $A(m)^2+C_H$.
This proves both the leading square and the full harmonic correction.

For the smooth part, the geometry of the cells supplies the correction.
Put $t=r/m$. The points $(\ell+\alpha)/n$, $0\leq\ell<r$,
are cell midpoints. We compute their cell endpoints explicitly. A cell centered at
$(\ell+\alpha)/n$ has endpoints
$(\ell+\alpha-1/2)/n$ and $(\ell+\alpha+1/2)/n$.
Since $\alpha=(1-t)/2$, the first left endpoint is $-t/(2n)$,
and the last right endpoint is
\[
 \frac{r+\alpha-1/2}{n}
 =\frac{r-t/2}{n}=t+\frac{t}{2n},
\]
using $r=mt=(n+1)t$. Their union is exactly
$[-t/(2n),t+t/(2n)]$.
On $[-1/4,5/4]$ the function $B$ and its first two derivatives are bounded.
The midpoint error in the sum is $O(r/n^2)=O(t/n)$.
The integrated extensions have length $t/(2n)$ on each side.
Taylor's formula for $B$ at $0$ and at $t$ gives
\[
 \begin{aligned}
 \int_{-t/(2n)}^0B(x)\,dx
    &=\frac{t}{2n}B(0)+O(t^2/n^2),\\
 \int_t^{t+t/(2n)}B(x)\,dx
    &=\frac{t}{2n}B(t)+O(t^2/n^2).
 \end{aligned}
\]
The midpoint error for the integral is $O(r/n^3)$, so its
error after multiplying by $n$ is $O(r/n^2)=O(t/n)$.
The extension errors after the same multiplication are
$O(t^2/n)\leq O(t/n)$. Combining all three estimates gives
\[
 \sum_{\ell<r}B((\ell+\alpha)/n)
   =nI(t)+\tfrac t2\{B(t)+B(0)\}+O(t/n),\qquad I(t)=\int_0^tB(x)\,dx.
\]
Put $W(t)=tw(t)$, smoothly extended by $W(0)=a$.
The functions $F_B(t)=w(t)I(t)=W(t)\int_0^1B(st)\,ds$ and
$G_B(t)=W(t)\{B(t)+1/\pi\}$ are smooth on $[0,1]$.
Their endpoint values are $F_B(0)=a^2/2$, $G_B(0)=a^2$,
$F_B(1)=G_B(1)=0$. Weighting the row expansion, using boundedness
of $W$, and applying smooth quadrature yields
\[
 \begin{aligned}
 \frac{4V_n}{nm}
 &=\frac4m\sum_{r<m}F_B(r/m)+\frac2{nm}\sum_{r<m}G_B(r/m)+O(n^{-2})\\
 &=4\int_0^1F_B-\frac{a^2}{m}
       +\frac2n\int_0^1G_B-\frac{a^2}{nm}+O(n^{-2}).
 \end{aligned}
\]
Since $1/m+1/(nm)=1/n$ and $\int G_B=I_B+I_2/\pi$, this is
\eqref{cu:V}, with its stated constants. The calculation keeps the last
row as well as both outer endpoint corrections.

The normalization in this final smooth calculation is worth
following once. The leading row term $nI(t)$ becomes
$4m^{-1}\sum F_B(t)$ after multiplication by $4/(nm)$.
Its quadrature endpoint correction is
$-2(F_B(0)+F_B(1))/m=-a^2/m$.
The extension term $t(B(t)+B(0))/2$ becomes
$2(nm)^{-1}\sum G_B(t)$, whose leading integral is
$2n^{-1}\int G_B$ and whose endpoint correction is
$-a^2/(nm)$.

The residual $O(t/n)$ in each row contributes at most
\[
 \frac{C}{n^2m}\sum_{r<m}w(r/m)(r/m)=O(n^{-2}),
\]
because $tw(t)$ is bounded and there are $m-1$ rows.
Thus all stated error estimates have the same constant across the
whole row range. The identity
\[
 \frac1m+\frac1{nm}=\frac{n+1}{nm}=\frac1n
\]
then combines the two endpoint corrections to $-a^2/n$.
Together with $\int G_B=I_B+I_2/\pi$, it gives
$\beta_B=-a^2+2I_B+aI_2$. This explains how a completely
smooth remainder still supplies a nonzero reciprocal correction.
\end{proof}

\begin{proof}[Proof of Theorem~\ref{thm:main-expansion}]
We now combine the two aggregate formulas with the exact scalar
$-1+2/m$. At this stage the objects $I_1,I_2,I_B$ are finite
ordinary integrals, so their algebraic combination is legitimate.
It is helpful to combine their integrands before integrating:
individual intermediate terms such as $w$ and $1/x$ would have
divergent integrals if separated. The displayed identity evaluates
a combination in which those endpoint divergences cancel.

Only one integral identity is needed to combine the two corrections:
\begin{equation}\label{cu:cancel}
 2I_B-aI_1=a^2-1+a^2\log(\pi/2).
\end{equation}
The double-angle identity $\cot(\pi x)=(w-w^{-1})/2$ makes its
combined integrand $x(w^2-1)-2aw+a^2/x$.
Since $w'=-(1+w^2)/a$, this equals
$-a(xw)'-a(w-a/x)-2x$. Integration gives \eqref{cu:cancel}, using
$\lim_{x\downarrow0}xw(x)=a$, $w(1)=0$ and
$\int_0^1(w-a/x)=a\log(2/\pi)$.

To verify the combined integrand, write
\[
 B(x)=\frac{w(x)-w(x)^{-1}}2-\frac{a}{2x}
                                      +\frac{a}{2(1-x)}.
\]
Multiplication by $2xw$ and subtraction of
$a\{w/(1-x)-a/x\}$ gives
\[
 x(w^2-1)-aw+\frac{a(x-1)w}{1-x}+\frac{a^2}{x}
     =x(w^2-1)-2aw+\frac{a^2}{x}.
\]
Differentiating cotangent gives $aw'=-(1+w^2)$.
Therefore
\[
 -a(xw)'-a(w-a/x)-2x
 =-2aw+x(1+w^2)+a^2/x-2x,
\]
which is the same expression. Integrate over $[\epsilon,1]$
and then let $\epsilon$ decrease to zero. The three resulting
terms are respectively
\[
 a^2,\qquad -a^2\log(2/\pi),\qquad -1.
\]
Their sum is $a^2-1+a^2\log(\pi/2)$, proving
\eqref{cu:cancel} while retaining the endpoint boundary value.

Replacing $\beta_B/n$ by $\beta_B/m$ costs $O(m^{-2})$.
The scalar $2/m$ in \eqref{cu:exact} now combines with the two
corrections as
\[
 2+\beta_H+\beta_B
 =-a^2-a^2(\gamma+2\log2)+2I_B-aI_1=-1-a^2d.
\]
Consequently, for a constant $C=-1+C_H+C_B$,
\[
 \Lambda_c(n)=A(m)^2+C-
       \frac{a^2\log m+1+a^2d}{m}+O(\log m/m^2).
\]
The limit \eqref{cu:kappa}, and $A(m)^2-A(n)^2\to0$, identify $C=\kappa$.

Here the roles of the two parts of the proof meet. The row
calculation has produced a degree-independent $C$ but has not
computed it. Subtract $A(n)^2$ from its result. Since
\[
 A(m)^2-A(n)^2
  =a^2\log(m/n)\{\log m+\log n+2d\}
                         =O(\log n/n),
\]
and the remaining reciprocal terms tend to zero, that difference
tends to $C$. The earlier independent mesh comparison proves that
the same difference tends to $\kappa$. Uniqueness of a real limit
therefore gives $C=\kappa$. This identification uses no numerical
approximation to the constant and does not assume the final
expansion.
The half-step variable $N=n+1/2$ is now forced by the logarithmic
reciprocal term. If $P(z)=A(z)^2$, then
\[
 P'(z)=\frac{2a^2(\log z+d)}{z},\qquad
 P''(z)=\frac{2a^2(1-\log z-d)}{z^2}.
\]
Taylor expansion from $m$ to $m-1/2$ subtracts
$P'(m)/2=(a^2\log m+a^2d)/m$ from the square. Its remainder
is $O(\log m/m^2)$ because $|P''(z)|$ has that bound on
$[m-1/2,m]$. Thus the shifted square absorbs exactly the
logarithmic reciprocal term and its associated constant. Applying
this calculation gives
\[
 A(N)^2=A(m)^2-\frac{a^2\log m+a^2d}{m}+O(\log m/m^2),
 \qquad N^{-1}=m^{-1}+O(m^{-2}).
\]
This proves $\Lambda_c(n)=A(N)^2+\kappa-N^{-1}+O(\log N/N^2)$.
Expanding the square gives \eqref{eq:main-expansion} with precisely the
constants \eqref{eq:main-constants}. Every bound holds for all $n\geq2$,
so the remainder is uniform across both parities. The corner proof is
entirely analytic; its application to $\Lambda_n$ uses
Theorem~\ref{thm:main-location}.

To make the final coefficient visible, write the preceding result
as
\[
 \Lambda_c(n)=A(N)^2+\kappa-\frac1N
                   +O(\log N/N^2).
\]
Expanding $A(N)^2$ gives
\[
 a^2\log^2N+2a^2d\log N+a^2d^2.
\]
The definitions $c_1=ad$, $b=2ac_1$, and $c=c_1^2+\kappa$
therefore yield the constants stated in the theorem. The correction
$-1/N$ is what remains after combining the exact endpoint scalar,
the harmonic endpoint sequence, the two smooth quadrature endpoints,
and the half-step shift. None of those pieces can be ignored when
computing this coefficient.

All errors have been bounded after their relevant weighted sums.
The small-$r$ and small-$s$ rows were explicitly included, and no
step selected a parity subsequence. Because $N$ and $m$ differ by
a fixed half unit and $m\geq3$, the functions
$\log m/m^2$ and $\log N/N^2$ bound one another up to fixed
factors. Thus there is one finite remainder constant for every
integer $n\geq2$. In particular the theorem also implies the
scaled limit
\[
 N\{\Lambda_c(n)-A(N)^2-\kappa\}\longrightarrow-1,
\]
since $\log N/N$ tends to zero. This limit describes the first
correction analytically. Passing from this corner value to the
maximum over the square uses the separate location theorem.
\end{proof}
